% Options for packages loaded elsewhere
\PassOptionsToPackage{unicode}{hyperref}
\PassOptionsToPackage{hyphens}{url}
\PassOptionsToPackage{dvipsnames,svgnames,x11names}{xcolor}
%
\documentclass[
  11pt,
  ignorenonframetext,
  aspectratio=169,
]{beamer}
\usepackage{pgfpages}
\setbeamertemplate{caption}[numbered]
\setbeamertemplate{caption label separator}{: }
\setbeamercolor{caption name}{fg=normal text.fg}
\beamertemplatenavigationsymbolsempty
% Prevent slide breaks in the middle of a paragraph
\widowpenalties 1 10000
\raggedbottom
\setbeamertemplate{part page}{
  \centering
  \begin{beamercolorbox}[sep=16pt,center]{part title}
    \usebeamerfont{part title}\insertpart\par
  \end{beamercolorbox}
}
\setbeamertemplate{section page}{
  \centering
  \begin{beamercolorbox}[sep=12pt,center]{part title}
    \usebeamerfont{section title}\insertsection\par
  \end{beamercolorbox}
}
\setbeamertemplate{subsection page}{
  \centering
  \begin{beamercolorbox}[sep=8pt,center]{part title}
    \usebeamerfont{subsection title}\insertsubsection\par
  \end{beamercolorbox}
}
\AtBeginPart{
  \frame{\partpage}
}
\AtBeginSection{
  \ifbibliography
  \else
    \frame{\sectionpage}
  \fi
}
\AtBeginSubsection{
  \frame{\subsectionpage}
}

\usepackage{amsmath,amssymb}
\usepackage{iftex}
\ifPDFTeX
  \usepackage[T1]{fontenc}
  \usepackage[utf8]{inputenc}
  \usepackage{textcomp} % provide euro and other symbols
\else % if luatex or xetex
  \usepackage{unicode-math}
  \defaultfontfeatures{Scale=MatchLowercase}
  \defaultfontfeatures[\rmfamily]{Ligatures=TeX,Scale=1}
\fi
\usepackage{lmodern}
\usetheme[]{Boadilla}
\usecolortheme{default}
\usefonttheme{serif} % use mainfont rather than sansfont for slide text
\ifPDFTeX\else  
    % xetex/luatex font selection
    \setmainfont[]{Latin Modern Sans}
\fi
% Use upquote if available, for straight quotes in verbatim environments
\IfFileExists{upquote.sty}{\usepackage{upquote}}{}
\IfFileExists{microtype.sty}{% use microtype if available
  \usepackage[]{microtype}
  \UseMicrotypeSet[protrusion]{basicmath} % disable protrusion for tt fonts
}{}
\makeatletter
\@ifundefined{KOMAClassName}{% if non-KOMA class
  \IfFileExists{parskip.sty}{%
    \usepackage{parskip}
  }{% else
    \setlength{\parindent}{0pt}
    \setlength{\parskip}{6pt plus 2pt minus 1pt}}
}{% if KOMA class
  \KOMAoptions{parskip=half}}
\makeatother
\usepackage{xcolor}
\newif\ifbibliography
\setlength{\emergencystretch}{3em} % prevent overfull lines
\setcounter{secnumdepth}{5}


\providecommand{\tightlist}{%
  \setlength{\itemsep}{0pt}\setlength{\parskip}{0pt}}\usepackage{longtable,booktabs,array}
\usepackage{calc} % for calculating minipage widths
\usepackage{caption}
% Make caption package work with longtable
\makeatletter
\def\fnum@table{\tablename~\thetable}
\makeatother
\usepackage{graphicx}
\makeatletter
\def\maxwidth{\ifdim\Gin@nat@width>\linewidth\linewidth\else\Gin@nat@width\fi}
\def\maxheight{\ifdim\Gin@nat@height>\textheight\textheight\else\Gin@nat@height\fi}
\makeatother
% Scale images if necessary, so that they will not overflow the page
% margins by default, and it is still possible to overwrite the defaults
% using explicit options in \includegraphics[width, height, ...]{}
\setkeys{Gin}{width=\maxwidth,height=\maxheight,keepaspectratio}
% Set default figure placement to htbp
\makeatletter
\def\fps@figure{htbp}
\makeatother

\makeatletter
\@ifpackageloaded{caption}{}{\usepackage{caption}}
\AtBeginDocument{%
\ifdefined\contentsname
  \renewcommand*\contentsname{Table of contents}
\else
  \newcommand\contentsname{Table of contents}
\fi
\ifdefined\listfigurename
  \renewcommand*\listfigurename{List of Figures}
\else
  \newcommand\listfigurename{List of Figures}
\fi
\ifdefined\listtablename
  \renewcommand*\listtablename{List of Tables}
\else
  \newcommand\listtablename{List of Tables}
\fi
\ifdefined\figurename
  \renewcommand*\figurename{Figure}
\else
  \newcommand\figurename{Figure}
\fi
\ifdefined\tablename
  \renewcommand*\tablename{Table}
\else
  \newcommand\tablename{Table}
\fi
}
\@ifpackageloaded{float}{}{\usepackage{float}}
\floatstyle{ruled}
\@ifundefined{c@chapter}{\newfloat{codelisting}{h}{lop}}{\newfloat{codelisting}{h}{lop}[chapter]}
\floatname{codelisting}{Listing}
\newcommand*\listoflistings{\listof{codelisting}{List of Listings}}
\makeatother
\makeatletter
\makeatother
\makeatletter
\@ifpackageloaded{caption}{}{\usepackage{caption}}
\@ifpackageloaded{subcaption}{}{\usepackage{subcaption}}
\makeatother

\ifLuaTeX
  \usepackage{selnolig}  % disable illegal ligatures
\fi
\usepackage{bookmark}

\IfFileExists{xurl.sty}{\usepackage{xurl}}{} % add URL line breaks if available
\urlstyle{same} % disable monospaced font for URLs
\hypersetup{
  pdftitle={Economía Política},
  pdfauthor={Sebastián Freille sfreille@unc.edu.ar Licenciatura en Economía FCE--UNC},
  colorlinks=true,
  linkcolor={blue},
  filecolor={Maroon},
  citecolor={Blue},
  urlcolor={Blue},
  pdfcreator={LaTeX via pandoc}}


\title{Economía Política}
\subtitle{Unidad}
\author{}
\date{}

\begin{document}
\frame{\titlepage}

\renewcommand*\contentsname{Contenidos}
\begin{frame}[allowframebreaks]
  \frametitle{Contenidos}
  \tableofcontents[hideallsubsections]
\end{frame}

\section{\texorpdfstring{1. \textbf{Polarización y Redes
Sociales}}{1. Polarización y Redes Sociales}}\label{polarizaciuxf3n-y-redes-sociales}

\begin{frame}{1. \textbf{Polarización y Redes Sociales}}
DOES SOCIAL MEDIA CAUSE POLARIZATION? EVIDENCE FROM ACCESS TO TWITTER
ECHO CHAMBERS DURING THE 2019 ARGENTINE PRESIDENTIAL DEBATE (Di Tella,
Gálvez, Schagrodski. 2021)
\end{frame}

\begin{frame}{1.1 Motivación:}
\phantomsection\label{motivaciuxf3n}
\begin{block}{1.1.1 ¿Cual fue el papel de las redes sociales en el
fomento de la polarización durante el debate presidencial de 2019 en
Argentina?}
\phantomsection\label{cual-fue-el-papel-de-las-redes-sociales-en-el-fomento-de-la-polarizaciuxf3n-durante-el-debate-presidencial-de-2019-en-argentina}
Hay una creciente preocupación acerca de que las personas vivimos dentro
de ``Burbujas partidistas'', con escasa exposición a quienes tienen una
inclinación ideológica diferente.

La polarización podría tener consecuencias negativas para la democracia,
el aumento de prejuicios y estancamiento legislativo.

~

Culpables \(\longrightarrow\) \textbf{Redes sociales??}
\end{block}
\end{frame}

\begin{frame}{1.2 Redes Sociales}
\phantomsection\label{redes-sociales}
Diversas investigaciones señalan que plataformas como \textbf{Facebook}
o \textbf{Twitter} tienden a generar \textbf{``cámaras de eco''},
espacios donde los usuarios están expuestos principalmente a
información, opiniones e ideas que refuerzan sus creencias y
perspectivas previas.

En estos entornos digitales, el intercambio de puntos de vista distintos
se ve limitado, lo que puede reducir la diversidad informativa y
favorecer la polarización de las opiniones.
\end{frame}

\begin{frame}{1.3 Experimento de Campo}
\phantomsection\label{experimento-de-campo}
\(\implies\) \textbf{Estudiar el papel de las redes sociales en el
fomento de la polarización en presencia de un evento políticamente
relevante: el debate presidencial argentino de 2019}

Se han propuesto dos estrategias para reducir el impacto aparente de las
redes sociales en la polarización:

\begin{enumerate}
[a.]
\tightlist
\item
  limitar el acceso de las personas a las cámaras de eco
\item
  garantizar que a través de las cámaras de eco reciban información que
  de otro modo no recibirían (datos ``contra-actitudinales'')
\end{enumerate}
\end{frame}

\begin{frame}{1.4 Experimento de Campo (Cont.)}
\phantomsection\label{experimento-de-campo-cont.}
\begin{quote}
Se centran en la forma en que una persona que ha formado una cámara de
eco política en Twitter en el período previo a las elecciones reacciona
a su ausencia en el momento preciso en que presumiblemente es más
valiosa
\end{quote}

\emph{Experimento:} ``el día del debate presidencial, se invitó a un
grupo de sujetos a acudir a un lugar neutral y participar en un estudio
académico a cambio de una compensación económica''.
\end{frame}

\begin{frame}{1.5 Experimento de Campo (Cont.)}
\phantomsection\label{experimento-de-campo-cont.-1}
\begin{block}{1.5.1 Pasos:}
\phantomsection\label{pasos}
\begin{enumerate}
\tightlist
\item
  Los participantes debieron completar un breve cuestionario y
  proporcionar una muestra de saliva
\item
  Fueron asignados aleatoriamente a cuatro grupos \(\implies\)
  \textbf{Tratamiento}
\item
  Los participantes debieron responder preguntas sobre polarización y el
  interés en la política y proporcionar una segunda muestra de saliva.
\end{enumerate}
\end{block}
\end{frame}

\begin{frame}{1.6 Tipos de participantes:}
\phantomsection\label{tipos-de-participantes}
\begin{itemize}
\tightlist
\item
  Usuarios de Twitter

  \begin{itemize}
  \tightlist
  \item
    Segregados
  \item
    No segregados
  \end{itemize}
\item
  No usuarios de Twitter(asignados a grupos 1 y 2)

  \begin{itemize}
  \tightlist
  \item
    Segregados
  \item
    No segregados
  \end{itemize}
\end{itemize}
\end{frame}

\begin{frame}{1.7 Segregación Usuarios Twitter}
\phantomsection\label{segregaciuxf3n-usuarios-twitter}
Para los participantes de Twitter, Se estima su preferencias de partido
mediante el análisis de sus interacciones con esta red y combinando la
información recopilada con las respuestas del cuestionario previo al
tratamiento se obtiene la segregación de los individuos.

\begin{figure}[H]

{\centering \includegraphics{../imagen/2.PNG}

}

\caption{Notas. La cifra incluye las 95 cuentas semilla y las 132.503
cuentas no semilla a las que siguen. Una escala de colores continua
refleja la orientación política prevista por los usuarios. El azul claro
(amarillo) indica una cuenta totalmente inclinada hacia FdT (JxC). El
tamaño de cada nodo refleja su grado de entrada.}

\end{figure}%
\end{frame}

\begin{frame}{1.8 Segregación de Usuarios No Tuiteros}
\phantomsection\label{segregaciuxf3n-de-usuarios-no-tuiteros}
El enfoque metodológico busca identificar, dentro del grupo de
\textbf{usuarios no tuiteros}, a aquellos individuos que, según sus
respuestas previas al tratamiento, \textbf{presentarían una tendencia a
segregarse} en caso de incorporarse a Twitter.

Para ello, se emplea un \textbf{modelo de aprendizaje supervisado}
estimado a partir de los datos observados en los usuarios activos de
Twitter, el cual permite predecir comportamientos de segregación en el
grupo no tuitero.
\end{frame}

\begin{frame}{1.9 Segregación de participantes}
\phantomsection\label{segregaciuxf3n-de-participantes}
\includegraphics{../imagen/1.PNG}
\end{frame}

\begin{frame}{1.10 Asignación de grupos y TRATAMIENTO}
\phantomsection\label{asignaciuxf3n-de-grupos-y-tratamiento}
Los participantes fueron divididos en cuatro grupos:

\begin{enumerate}
\item
  \textbf{Documentary group}: Apagaron sus celulares y vieron un corto
  documental sobre las instituciones políticas argentinas.
\item
  \textbf{Debate-Only}: Apagaron sus celulares, pero vieron la
  transmisión en vivo del debate presidencial.
\item
  \textbf{Twitter-Allowed}: Pudieron ver el debate en condiciones
  normales, usando sus celulares libremente y accediendo a Twitter.
\item
  \textbf{Twitter-Interact}: Pudieron ver el debate y utilizar sus
  celulares, pero en este caso, se les pidió que interactuaran con una
  ``cámara de eco'' que compartía contenido polarizado
  (``contra-actitudinales'').
\end{enumerate}
\end{frame}

\begin{frame}{1.11 Segunda Encuesta}
\phantomsection\label{segunda-encuesta}
Después de los tratamientos, los sujetos tuvieron que responder a una
serie de preguntas estándar sobre la polarización y el interés en la
política, así como proporcionar una segunda muestra de saliva.
\end{frame}

\begin{frame}{1.12 Indices de polarización:}
\phantomsection\label{indices-de-polarizaciuxf3n}
\begin{enumerate}
\item
  \textbf{Índice Élite:} Desagrado que los participantes sienten hacia
  las élites del partido externo. \emph{Dos componentes principales:
  Percepción de amenaza y Comparación de rasgos positivos}.
\item
  \textbf{Índice de Distancia Social:} Desagrado que los ciudadanos
  sienten hacia los partidarios del partido opuesto. \emph{Tres
  componentes: Familia, Socialización, Colaboración laboral}.
\item
  \textbf{Índice de Partidismo:} La convicción que tienen los
  participantes sobre su elección política. \emph{Dos componentes:
  Convicción del voto, Dispersión de preferencias}.
\end{enumerate}
\end{frame}

\begin{frame}{1.13 Indices de polarización (Cont.):}
\phantomsection\label{indices-de-polarizaciuxf3n-cont.}
\begin{enumerate}
\setcounter{enumi}{3}
\item
  \textbf{Índice Institucional:} Desconfianza que los participantes
  sienten hacia las instituciones. \emph{Tres componentes: Credibilidad
  de acusaciones de fraude, Desacuerdo con sentencias judiciales,
  Lealtad vs.~competencia}.
\item
  \textbf{Índice de Compromiso:} Grado de disposición de los
  participantes a involucrarse en el proceso electoral. \emph{Dos
  componentes: Participación activa,Persuasión}.
\item
  \textbf{Índices Combinados:}Un índice de polarización que combina
  diferentes dimensiones.
\end{enumerate}
\end{frame}

\section{2. Resultados}\label{resultados}

\begin{frame}{2.1 Fuentes de Información}
\phantomsection\label{fuentes-de-informaciuxf3n}
Hay 3 tipos de datos obtenidos:

\begin{itemize}
\tightlist
\item
  Resultados estándar basados en encuestas.
\item
  Medida de los niveles de cortisol.\\
\item
  Y para los usuarios de Twitter, información sobre su actividad en
  Twitter.
\end{itemize}
\end{frame}

\begin{frame}{2.2 Polarización}
\phantomsection\label{polarizaciuxf3n}
\begin{figure}[H]

{\centering \includegraphics{../imagen/3.PNG}

}

\caption{Notas:Todas las variables se normalizan para que la
distribución de grupos permitidos por Twitter no segregada tenga una
media de 0 y una desviación estándar de 1. Las barras de error reflejan
intervalos de confianza del 90\%}

\end{figure}%
\end{frame}

\begin{frame}{2.3 Efectos del tratamiento}
\phantomsection\label{efectos-del-tratamiento}
\includegraphics{../imagen/4.PNG}
\end{frame}

\begin{frame}{2.4 Efectos del tratamiento}
\phantomsection\label{efectos-del-tratamiento-1}
\begin{itemize}
\tightlist
\item
  Para usuarios de Twitter no segregados, no hay diferencias en el
  índice de polarización combinado para ningún tratamiento.
\item
  Para de usuarios segregados de Twitter, observamos que el
  \textbf{índice de polarización combinado más bajo que medimos se
  encuentra en el tratamiento Twitter-Permitido.}

  \begin{itemize}
  \tightlist
  \item
    aunque la diferencia con el del grupo Documental no es
    estadísticamente significativa.
  \item
    la diferencia en los niveles de polarización entre Twitter-Permitido
    y los otros dos tratamientos (Twitter-Solo Debate y
    Twitter-Interact) son cómodamente significativas
  \end{itemize}
\item
  Para no usuarios de Twitter, \textbf{los niveles de polarización
  agregados son similares en los dos tratamientos} tanto en la muestra
  de usuarios segregados como en la no segregada.
\end{itemize}
\end{frame}

\begin{frame}{2.5 Conclusiones}
\phantomsection\label{conclusiones}
\begin{itemize}
\tightlist
\item
  El efecto del debate, aproximado mediante la comparación de los dos
  tratamientos disponibles (Solo Debate vs Documental) sugiere que el
  evento no cambia los niveles de polarización. Curiosamente, tampoco lo
  hace con los usuarios de Twitter.
\item
  Para no usuarios de Twitter el \textbf{grupo segregado está
  significativamente más polarizado que el no segregado}, lo que sugiere
  que la polarización no es exclusivamente un fenómeno en línea.
\item
  Parece que, en promedio, hay una polarización significativamente menor
  en los no Twitter en relación con la muestra de Twitter. (Los usuarios
  de Twitter son mas ``politizados''?).
\item
  Para sujetos que ya estaban en cámaras de eco muestran una mayor
  polarización para el grupo de Solo Debate como para el de
  Twitter-Interact.
\end{itemize}
\end{frame}

\begin{frame}{2.6 Conclusiones}
\phantomsection\label{conclusiones-1}
\begin{itemize}
\tightlist
\item
  Para el grupo de Twitter-Interact, también hay un aumento en ΔCortisol
  (lo que sugiere que las tareas involucradas fueron una fuente de
  estrés para estos participantes) y en la segregación de ΔContact (lo
  que sugiere que sus interacciones en Twitter se volvieron más
  polarizadas).
\end{itemize}

\textbf{Nuestro principal resultado puede interpretarse como la
identificación de circunstancias en las que la presencia de una cámara
de eco puede verse como despolarizante. ¿A qué se debe? El mecanismo que
imaginamos es que las personas que están interesadas en la política
``demanden'' cámaras de eco porque les es útil cuando tienen que
interpretar un evento político. }

\begin{itemize}
\tightlist
\item
  En todo caso, hay alguna evidencia débil de una reducción en la
  polarización de sus contactos de Twitter en la condición
  contra-actitudinal.
\item
  Nuestra conjetura es que las cámaras de eco de Twitter que los sujetos
  han formado con otros usuarios de Twitter de ideas afines les son
  útiles cuando tienen que interpretar un evento político
\end{itemize}
\end{frame}




\end{document}
